\documentclass[11pt,a4paper,openany]{ctexbook}
\usepackage[a4paper,margin=1.8cm,headheight=16pt]{geometry}
\usepackage{amsmath,amssymb,mathtools}
\usepackage{booktabs,longtable,tabularx,array,multirow}
\newcolumntype{L}[1]{>{\raggedright\arraybackslash}p{#1}}
\setlength{\extrarowheight}{2pt}
\usepackage{enumitem}
\usepackage{tcolorbox}
\tcbuselibrary{breakable,skins,theorems}
\usepackage{tikz}
\usetikzlibrary{arrows.meta,positioning,calc,shapes.geometric,fit,matrix}
\usepackage{xcolor}
\usepackage{fancyhdr}
\usepackage{hyperref}
\usepackage{microtype}
\usepackage{titlesec}
\usepackage{lastpage}

\definecolor{ink}{RGB}{25,25,25}
\definecolor{deep}{RGB}{42,58,73}
\definecolor{soft}{RGB}{245,247,249}
\definecolor{line}{RGB}{150,158,166}
\definecolor{accent}{RGB}{104,76,55}
\definecolor{warn}{RGB}{123,76,67}
\definecolor{greenish}{RGB}{57,92,76}
\definecolor{blueish}{RGB}{68,92,126}

\hypersetup{colorlinks=true,linkcolor=deep,urlcolor=accent,citecolor=deep,
  pdftitle={408统一总图：四科心智模型、章节映射与做题入口}}
\setlength{\parindent}{2em}
\setlength{\parskip}{0.25em}
\setlength{\emergencystretch}{3em}
\renewcommand{\arraystretch}{1.25}
\setlength{\tabcolsep}{4.5pt}
\setlist{itemsep=0.15em,topsep=0.35em,leftmargin=2em}
\pagestyle{fancy}
\fancyhf{}
\lhead{\small\color{deep}\textbf{408 统一总图}}
\rhead{\small\color{deep}\leftmark}
\cfoot{\small 第 \thepage\ 页}
\renewcommand{\headrulewidth}{0.4pt}

\titleformat{\chapter}[display]
  {\normalfont\huge\bfseries\color{deep}}
  {\filleft\Large 第\thechapter 篇}
  {0.4ex}
  {\titlerule\vspace{0.8ex}\filright}
\titleformat{\section}{\Large\bfseries\color{deep}}{\thesection}{0.6em}{}
\titleformat{\subsection}{\large\bfseries\color{ink}}{\thesubsection}{0.6em}{}

\newtcolorbox{corebox}[1]{breakable,colback=soft,colframe=deep,title=\textbf{#1},boxrule=0.8pt,arc=2mm,left=2mm,right=2mm,top=1mm,bottom=1mm}
\newtcolorbox{methodbox}[1]{breakable,colback=white,colframe=greenish,title=\textbf{#1},boxrule=0.7pt,arc=1.5mm}
\newtcolorbox{warnbox}[1]{breakable,colback=white,colframe=warn,title=\textbf{#1},boxrule=0.7pt,arc=1.5mm}
\newtcolorbox{examBox}[1]{breakable,colback=white,colframe=accent,title=\textbf{408 训练：#1},boxrule=0.7pt,arc=1.5mm}
\newtcolorbox{boundarybox}[1]{breakable,colback=soft,colframe=line,title=\textbf{总图 / 专题边界：#1},boxrule=0.6pt,arc=1.5mm}

\newcommand{\key}[1]{\textbf{\color{deep}#1}}
\newcommand{\eng}[1]{\textit{#1}}
\newcommand{\arrowchain}[1]{\begin{center}\fbox{\parbox{0.92\textwidth}{\centering\ttfamily #1}}\end{center}}
\let\cleardoublepage\clearpage

\tikzset{
  atlas/.style={draw=deep,rounded corners=2mm,minimum width=3.25cm,minimum height=1.0cm,align=center,fill=soft,font=\small},
  bridge/.style={draw=greenish,rounded corners=2mm,minimum width=3.25cm,minimum height=0.95cm,align=center,fill=white,font=\small},
  flow/.style={-{Latex[length=2mm]},thick,draw=deep}
}

\begin{document}
\frontmatter

\begin{titlepage}
\centering
\vspace*{1.2cm}
{\Huge\bfseries\color{deep} 408 统一总图\par}
\vspace{0.45cm}
{\LARGE 四科心智模型 · 章节映射 · 做题入口\par}
\vspace{0.6cm}
{\large 数据结构 · 计算机组成原理 · 操作系统 · 计算机网络\par}
\vspace{1.4cm}
\begin{tikzpicture}[x=4.0cm,y=1.65cm]
  \node[atlas] (ds) at (0,0) {数据结构\\\small 组织信息};
  \node[atlas] (co) at (1,0) {计算机组成\\\small 执行指令};
  \node[atlas] (os) at (2,0) {操作系统\\\small 管理机器};
  \node[atlas] (net) at (3,0) {计算机网络\\\small 跨机协作};
  \draw[flow] (ds) -- node[above,font=\scriptsize]{表示影响真实访问} (co);
  \draw[flow] (co) -- node[above,font=\scriptsize]{硬件机制} (os);
  \draw[flow] (os) -- node[above,font=\scriptsize]{socket / NIC} (net);
  \node[bridge] (u) at (1.5,-1.15) {统一观察语言\\对象 · 表示/状态 · 变化 · 不变量 · 代价};
  \draw[flow] (ds) -- (u);
  \draw[flow] (co) -- (u);
  \draw[flow] (os) -- (u);
  \draw[flow] (net) -- (u);
\end{tikzpicture}
\vfill
{\small 本册负责 Map：建立学科坐标系、概念边界与专题接口；专题册负责 Depth。\par}
\end{titlepage}

\chapter*{使用说明：这本总手册到底负责什么}
\addcontentsline{toc}{chapter}{使用说明：这本总手册到底负责什么}

\begin{corebox}{总手册的唯一任务}
本册不是四本教材的压缩版，也不是所有专题笔记的合集。它只负责回答四类问题：
\begin{enumerate}
  \item \key{这门学科究竟研究什么世界？}
  \item \key{这门学科最少需要保留哪几个生成性母模型？}
  \item \key{一个知识点属于哪个专题，和相邻专题在哪里交接？}
  \item \key{面对陌生综合题，第一次应该从什么坐标进入推理？}
\end{enumerate}
\end{corebox}

\begin{boundarybox}{四类手册各司其职}
\begin{longtable}{L{2.6cm}L{4.0cm}L{\dimexpr\linewidth-7.85cm\relax}}
\toprule
\textbf{类型} & \textbf{任务} & \textbf{不应该做什么}\\
\midrule
总图 Atlas & 建立母问题、母模型、专题地图、关键边界 & 不展开具体算法、协议步骤、寄存器级流程\\
专题 Topic & 打穿一个核心机制，解释对象、状态、过程、不变量、计算模型 & 不重新建立整门学科总图\\
桥梁 Bridge & 解释两个专题/学科的接口和状态交接 & 不重复两边全部知识\\
综合 Integration & 在一个真实过程里同时调用多个已学机制 & 不成为新的知识大全\\
\bottomrule
\end{longtable}
\end{boundarybox}

本册中的“完整”因此不是细节最多，而是\key{地图功能完整}：总定义、核心坐标、专题所有权、关键边界、跨科接口和做题入口全部齐备；详细机制与覆盖清单下沉到四科方法论及各专题手册。

\tableofcontents
\mainmatter

% ============================================================
\chapter{统一底座：四门课不是同一个模型，但可以用同一组问题观察}
% ============================================================

\section{不要强行用一条公式统一四科}
四门课确实共享“对象、状态、映射、不变量、成本”等思想，但它们研究的核心矛盾不同。真正稳妥的统一方式不是要求所有知识都套进同一个状态机，而是建立一个\key{观察镜头}：先用少量问题确定自己正在研究什么，再切换到该学科自己的母模型。

\begin{corebox}{408 通用五问：这是观察顺序，不是万能公式}
面对任何新概念，先问：
\begin{enumerate}
  \item \key{Object / Goal：}现在研究的对象与最终目标是什么？
  \item \key{Representation / State：}为了完成目标，系统保存了哪些信息？这些信息存在哪里？
  \item \key{Operation / Event：}什么操作或事件会改变这些信息？
  \item \key{Invariant / Boundary：}什么条件必须始终成立，哪些相似概念绝不能混用？
  \item \key{Cost / Workload：}时间、空间、带宽、I/O、等待或可靠性代价来自哪里？
\end{enumerate}
\end{corebox}

这里的五问不是要求每道题都机械写五行答案，而是防止一上来就套公式。例如做“TLB + 页表 + Cache”综合题时，第一问先确认对象是一次虚拟地址访问；第二问确认题目给出的状态包括 TLB、页表和 Cache；第三问确认事件是 CPU 发出一次访存；第四问区分 TLB miss、page fault、cache miss；第五问最后才计算访问次数或有效访问时间。换句话说，\key{先确定世界，再进入计算}。

\section{本册符号到底怎样读}
为了压缩复杂关系，本册会使用箭头和英文关键词。但它们只是一种\key{思考记号}，不是要求背诵的形式语言。第一次看到时按下面方式理解。

\begin{longtable}{L{3.0cm}L{5.1cm}L{\dimexpr\linewidth-9.37cm\relax}}
\toprule
\textbf{记号} & \textbf{在本册中的含义} & \textbf{怎样实际使用}\\
\midrule
$A\rightarrow B$ & 从 $A$ 进入 $B$ 的\key{推理顺序、转换关系或依赖关系}；具体含义由上下文决定，不等于“物理上 A 一定先发生” & 例如 $Goal\rightarrow Representation$ 表示先知道目标，再检查表示；$Name\rightarrow IP$ 则表示名字经过解析得到网络地址 \\
\midrule
$A\xrightarrow{E}B$ & 当前状态 $A$ 在事件/操作 $E$ 发生后变成 $B$ & 例如 $Blocked\xrightarrow{I/O\ completion}Ready$；做状态题时必须能说出箭头上是什么事件 \\
\midrule
$A\neq B$ & 两个概念\key{可能表面相似，但判定问题、状态所有者或作用域不同，不能互换} & 看到题目时要找“决定它属于 A 还是 B 的判据”，而不是只背一句不等式 \\
\midrule
$A+B+C$ & 三个因素需要\key{共同存在或共同考虑}，不是数值加法 & 如 Seq+ACK+Timer 表示可靠传输至少同时依赖这些状态/机制 \\
\midrule
$A/B$ & 常用于并列两个观察维度或术语，表示“同时注意 A 与 B”，不是除法 & 如 Mechanism/Policy 表示分析时分别判断“能怎样做”和“选择怎样做” \\
\midrule
$A\leftrightarrow B$ & 两侧存在接口、双向影响或反馈 & 如 OS $\leftrightarrow$ Network 表示 socket、buffer、NIC 等处存在状态交接 \\
\bottomrule
\end{longtable}

\begin{methodbox}{读箭头的固定方法}
每看到一条母链，不要朗读英文单词。把每个框翻译成一句问题：\textbf{“这一格让我现在看什么？”}；把每个箭头翻译成：\textbf{“确认完前一格以后，下一步为什么有资格看后一格？”}。如果说不出这两句话，就说明这条模型还没有真正理解。
\end{methodbox}

\section{四门课各自真正拥有的母模型}

\begin{corebox}{数据结构：为操作需求选择可维护的表示}
\[
\boxed{\text{Problem}\rightarrow\text{Relation}\rightarrow\text{Operations}\rightarrow\text{Representation}\rightarrow\text{Invariant}\rightarrow\text{Cost}}
\]
这条链的中文意思是：\key{先明确题目要解决什么，再看数据之间是什么关系；由需要频繁执行的操作决定应该怎样表示这些关系；算法每次修改表示时必须维护结构不变量；最后在具体 workload 下评价成本。}
\end{corebox}

例如“需要不断插入元素并反复取当前最大值”：目标不是把所有元素完全排序，而是支持 $Insert+FindMax/DeleteMax$；因此只维护完全二叉树上的堆偏序即可，用数组保存；$FindMax$ 为常数时间，插入/删除极值通过上滤/下滤恢复 invariant。这个例子会分别在“堆”“数组表示”“复杂度”专题继续展开。

\begin{corebox}{计算机组成原理：把 ISA 语义实现成真实硬件时序}
\[
\boxed{\text{ISA State}\rightarrow\text{Data Location}\rightarrow\text{Datapath/Control}\rightarrow\text{Timing}\rightarrow\text{Commit}\rightarrow\text{Cost}}
\]
中文意思是：\key{先确定指令对软件承诺改变什么状态，再找操作数当前在哪里；然后追踪硬件经过哪些数据通路、由哪些控制信号驱动；判断值在什么周期产生和可用；最后确认何时对 ISA 可见并计算性能代价。}
\end{corebox}

例如一条 \texttt{LOAD R1, 8(R2)}：软件语义是把地址 $R2+8$ 处的数据写入 $R1$；硬件必须先读 $R2$、算有效地址、访问 Cache/Memory、等待数据返回，再在允许的时钟边界写回 $R1$。流水线、Cache、数据冒险和 CPI 都只是这条路径不同位置的问题。

\begin{corebox}{操作系统：在事件中维护受保护的资源世界}
\[
\boxed{\text{Objects/Relations/Queues}\xrightarrow{\text{Event + Mechanism + Policy}}\text{New State}}
\]
中文意思是：\key{先画清系统里有哪些对象、它们怎样引用/映射、谁正在排队等待；事件发生以后，硬件/内核用某种机制完成合法状态转换，策略决定多个合法选择中选哪一个；转换前后都必须满足保护、正确性和生命周期约束。}
\end{corebox}

例如进程因 \texttt{read()} 等待磁盘：Process 是对象，等待队列记录“它在等谁”；I/O completion 是事件；wakeup 是机制，把它从 Blocked 变成 Ready；scheduler policy 之后才决定什么时候让它从 Ready 变成 Running。这样就不会把“唤醒”和“立刻运行”混成一件事。

\begin{corebox}{计算机网络：没有全局即时状态时的分布式协作}
\[
\boxed{\text{Scope}\rightarrow\text{Local State}\xrightarrow{\text{Message/Timer}}\text{Transition}\rightarrow\text{Feedback}\rightarrow\text{Cost}}
\]
中文意思是：\key{先确定当前决策负责多远，再确定做决策的设备手里有哪些局部状态；收到报文或定时器到期后，它只依据这些本地信息更新状态并采取动作；动作通过新的报文影响其他节点；最后分析时延、带宽、队列和可靠性代价。}
\end{corebox}

例如 TCP 超时重传：作用域是端到端；发送方保存未确认序号、定时器等本地状态；timer expiration 触发重传；新的 segment/ACK 构成反馈；性能上付出额外时延和带宽。这就是“局部状态 + 消息/定时器”的具体落地。

\section{母模型怎样对应后续章节}
母模型不是替代章节目录，而是告诉你\key{学习某一章时，到底在母模型的哪一格继续向下挖}。

\begin{longtable}{L{2.0cm}L{4.1cm}L{4.2cm}L{\dimexpr\linewidth-11.57cm\relax}}
\toprule
\textbf{学科} & \textbf{母模型中的位置} & \textbf{典型章节} & \textbf{学习时真正要问}\\
\midrule
DS & Relation / Representation & 线性表、树、图、散列、B+树 & 逻辑关系如何编码？保存了哪些额外信息？ \\
DS & Invariant / Cost & 排序、平衡树、图算法、复杂度 & 每一步怎样保持正确？为了快付出什么？ \\
CO & Data Location / Path & 寄存器、主存、Cache、数据通路 & 数据现在在哪，下一步经过谁？ \\
CO & Timing / Commit & 流水线、异常、中断、性能 & 值何时可用，何时算真正完成？ \\
OS & Objects / Relations & 进程、VM、文件系统 & 对象身份是什么，名字/映射怎样指向它？ \\
OS & Queues / Event / Policy & 调度、并发、I/O、页面置换 & 谁在等，什么事件改变资格，策略选谁？ \\
NET & Scope / State Owner & 分层、MAC/IP/Port、路由、TCP & 当前决定负责多远，状态在端点/交换机/路由器谁手里？ \\
NET & Message/Timer / Feedback & 可靠传输、路由收敛、拥塞 & 收到什么或等多久以后改变状态，怎样影响别的节点？ \\
\bottomrule
\end{longtable}

\section{四科在完整计算系统里的分工}
数据结构与另外三科不是同一条硬件“层级”。更准确的关系是：\key{数据结构提供组织状态的通用工具；计组、OS、网络分别赋予这些状态硬件、内核和分布式语义。}

\begin{center}
\begin{tikzpicture}[x=4.0cm,y=1.65cm]
  \node[atlas] (co) at (0,0) {计组\\\small 单机硬件如何执行};
  \node[atlas] (os) at (1,0) {OS\\\small 单机资源如何管理};
  \node[atlas] (net) at (2,0) {网络\\\small 多机如何通信};
  \draw[flow] (co) -- node[above,font=\scriptsize]{ISA/MMU/中断/DMA} (os);
  \draw[flow] (os) -- node[above,font=\scriptsize]{socket/driver/buffer} (net);
  \node[bridge] (ds) at (1,-1.1) {数据结构：队列、树、图、散列、索引\\\small 为三门系统课组织状态并提供操作成本};
  \draw[flow] (ds) -- (co);
  \draw[flow] (ds) -- (os);
  \draw[flow] (ds) -- (net);
\end{tikzpicture}
\end{center}

图中的箭头表示\key{知识接口}而不是“学习顺序”。例如数据结构中的 Queue 在 OS 里承载 Ready/Wait Queue，在网络里承载 packet buffer；队列的数据结构性质决定插入/取出成本，但“为什么谁应该先出队”属于 OS/网络自己的策略语义。

\section{跨四科反复出现的六种设计思想：必须落到实例}
总图保留六条横向规律，但每一条都必须能在具体章节找到实例，否则只是抽象口号。

\begin{longtable}{L{3.0cm}L{4.3cm}L{4.3cm}L{\dimexpr\linewidth-12.87cm\relax}}
\toprule
\textbf{设计思想} & \textbf{具体例子} & \textbf{为什么有用} & \textbf{对应专题}\\
\midrule
Representation / Indirection & Hash: Key$\to$Bucket；OS: VA$\to$PTE$\to$Frame；Net: Name$\to$IP & 多一层映射换来更快定位、重定位、共享或抽象隔离 & DS 散列；CO Cache；OS VM/FS；NET DNS/IP \\
\midrule
State / Invariant & Heap 父子偏序；流水线最终 ISA 语义；PV 缓冲区容量；TCP 序号窗口 & 解释“为什么某一步能做/不能做”，是验证答案的依据 & DS 算法；CO 流水线；OS 并发；NET 可靠传输 \\
\midrule
Scarcity / Sharing & ALU/总线冲突；CPU ready queue；共享链路 packet queue & 需求超过瞬时供给时必然出现等待、排队或仲裁 & CO hazard；OS 调度；NET MAC/拥塞 \\
\midrule
Locality / Caching & Cache line；TLB/Page Cache；DNS cache & 用历史/邻近访问换取平均更快，但要处理 miss、替换和一致性 & CO Cache；OS VM/IO；NET DNS \\
\midrule
Fast / Slow Path & Hash 无冲突/冲突；Cache hit/miss；TLB hit/page fault；正常 ACK/timeout & 常见情况走短路径，少见失败承担复杂修复成本 & 四科多个专题 \\
\midrule
Trade-off / Workload & Hash 点查快但不擅长范围；流水线吞吐高但有 hazard；RR 响应好但切换多；窗口大提高吞吐也增加在途状态 & “最好”必须指定工作负载和成本目标 & 四科性能/算法选择题 \\
\bottomrule
\end{longtable}

\section{总手册统一做题入口：箭头是检查顺序}
总手册只负责\key{第一轮定位}，进入专题以后使用专题自己的解题协议。下面的箭头表示“前一项确认以后，再进入后一项”，不是说题目中的物理过程一定按这个顺序发生。

\begin{examBox}{跨科题第一次定位七问}
\[
\boxed{\begin{gathered}
Target\rightarrow Scope/Topic\rightarrow GivenState\rightarrow Event/Operation\\
\downarrow\\
Invariant/Boundary\rightarrow Mechanism/Path\rightarrow Cost
\end{gathered}}
\]
\begin{enumerate}
  \item \key{Target：}题目最终要你求数值、判断状态，还是选择结构/机制？先把问号圈出来。
  \item \key{Scope/Topic：}它属于哪门课、哪个专题、哪个层次？例如“下一跳”先进入网络转发，不要先想 TCP。
  \item \key{GivenState：}题目已经给了哪些表、队列、寄存器、指针、窗口、Cache 行？先抄出真正会变化的状态。
  \item \key{Event/Operation：}是谁做了什么，或者什么事件发生？没有事件依据，不要擅自改变状态。
  \item \key{Invariant/Boundary：}有哪些规则限制答案？例如 heap 偏序、页内偏移不变、Ready/Blocked 不可混。
  \item \key{Mechanism/Path：}沿哪条表示、映射或数据通路执行？需要时画表、时间线或状态图。
  \item \key{Cost：}最后才把路径转换成比较次数、cycle、I/O、时延、吞吐等数字。
\end{enumerate}
\end{examBox}

\begin{methodbox}{小例子：为什么先定位再计算}
题目问“某虚拟地址访问时 TLB 未命中，页表项有效，Cache 命中，主存访问几次？”先按入口链：Target=访存次数；Topic=地址翻译+Cache；GivenState=TLB miss/PTE valid/Cache hit；Event=一次 load；Boundary=TLB miss 不等于 page fault；Path=查页表得到 PA，再访问 Cache；最后才根据题设的硬件访问模型计数。若第一步就把 TLB miss 当成 page fault，后面的数字算得再熟也会整题走错。
\end{methodbox}

% ============================================================
\chapter{数据结构：为了操作而保存信息}
% ============================================================

\section{学科母问题}
\begin{corebox}{数据结构真正研究什么}
给定一组对象、它们之间的逻辑关系以及需要频繁完成的操作，\key{应该保存哪些信息、采用什么表示，使操作在维持结构语义的同时具有可接受成本？}
\end{corebox}

所以数据结构不是“容器目录”，而是：
\[
\boxed{\text{Problem}\rightarrow\text{Relation}\rightarrow\text{Operations}\rightarrow\text{Representation}\rightarrow\text{Invariant}\rightarrow\text{Cost}}
\]

\subsection{逐格解释这条母链}
\begin{description}
  \item[Problem] 题目真正要解决的任务，例如“频繁找最大值”“判断连通”“支持范围查询”。
  \item[Relation] 数据之间客观需要表达的关系，例如线性前后关系、树的层次关系、图的任意连接、集合归属关系。
  \item[Operations] 使用者真正频繁执行的操作集合以及频率。它决定哪些能力值得提前维护。
  \item[Representation] 计算机怎样把关系编码进数组位置、指针、矩阵、散列桶或索引节点。
  \item[Invariant] 每次插入、删除、交换、旋转后仍必须成立的结构条件；它是判断算法是否正确的核心。
  \item[Cost] 为这些能力付出的比较、移动、访存、空间、递归深度或外存 I/O。
\end{description}
箭头表示\key{设计/解题时的因果顺序}：不能脱离目标操作先宣布“用链表”或“用 Hash”；先有需求，才有表示选择。

\section{一个例子把母模型完整跑一遍：为什么会有堆}
假设任务是：数据不断到来，需要随时知道当前最大值，并经常删除最大值。

\begin{longtable}{L{2.6cm}L{\dimexpr\linewidth-3.45cm\relax}}
\toprule
\textbf{母模型位置} & \textbf{这个例子里具体是什么}\\
\midrule
Problem & 动态维护一批元素，同时高频执行“取最大、插入、删除最大”。\\
Operations & $FindMax+Insert+DeleteMax$；并不要求任意元素之间完全有序。\\
Relation & 只需要维护“父节点不小于孩子”的局部偏序，并利用完全二叉树保证紧凑。\\
Representation & 完全二叉树可以直接用数组按下标表示父子位置，不必保存大量指针。\\
Invariant & 对任意父子，$Parent\ge Child$。插入/删除破坏后用上滤/下滤恢复。\\
Cost & $FindMax=O(1)$，插入和删除最大值沿树高调整，约 $O(\log n)$；代价是不能像完全有序数组那样直接做有序遍历。\\
\bottomrule
\end{longtable}

因此学习“堆”时，不应只背数组下标和上滤代码，而要看见：\key{它是 Operation Set 决定 Representation 与 Invariant 的一个标准实例。}

\section{四个生成性支点，以及它们在哪些章节出现}
\subsection{Relation 与 Representation 分离}
线性、层次、任意网络关系是逻辑结构；数组、链表、树节点、邻接矩阵/表是表示。\key{同一逻辑关系可有多种表示，表示决定实际可见操作与成本。}

实际章节：线性表里比较顺序表/链表；图里比较邻接矩阵/邻接表；树可以顺序或链式表示；外存索引则把“树”映射到磁盘页/块。

\subsection{Operation Set 决定结构价值}
堆、并查集、散列、B+ 树之所以不同，不是因为形状不同，而是它们主动支持的核心操作不同。问“哪个结构最好”之前必须先问 workload。

实际章节：Hash 强点查；B+ 树兼顾点查与范围并降低 I/O；并查集只保留集合代表关系；堆只维护极值所需偏序。

\subsection{Invariant 是算法的稳定内核}
BST 的有序性、Heap 的偏序、AVL 的平衡、并查集的代表元、图算法的已确定区域，都是算法真正维护的对象。代码只是实现不变量恢复的一种形式。

实际章节：快速排序的一轮 partition、Dijkstra 的已确定最短距离集合、Kruskal 的无环性、循环队列的空/满规则，都可以用“不变量”检查步骤是否合法。

\subsection{Cost 不是单个 $O(\cdot)$ 标签}
比较、移动、随机访问、额外空间、递归深度、cache locality、外存 I/O 都可能成为真正成本；当成本模型从 RAM 变成磁盘 I/O 时，最佳结构也会改变。

实际章节：内部排序主要关心比较/移动；外部排序关心 I/O 趟数；B/B+ 树关心树高和块访问；数组与链表即使渐近复杂度相同，机器局部性也可能不同。

\section{母模型—章节映射：学每一章时到底在深入哪一格}
\begin{longtable}{L{3.0cm}L{4.2cm}L{\dimexpr\linewidth-8.47cm\relax}}
\toprule
\textbf{章节/专题} & \textbf{主要对应母模型位置} & \textbf{学习时抓住的问题}\\
\midrule
复杂度 & Cost & 输入规模是什么？基本操作是什么？最好/平均/最坏/摊还口径是否一致？\\
线性表/数组 & Relation $\to$ Representation & 相同线性关系为什么连续表示和链式表示的定位/插删代价不同？\\
栈/队列 & Operations $\to$ Semantics & 为什么限制操作端点反而得到 LIFO/FIFO 的过程语义？\\
树/二叉树 & Relation + Invariant & 层次关系怎样带来递归结构，不同树靠什么 invariant 区分？\\
图遍历 & State/Progress & Visited、Frontier、ExpansionOrder 如何共同推进探索？\\
图算法 & Invariant & 哪些结论已永久确定，局部松弛/选边为什么不会破坏全局正确性？\\
查找/散列/B+树 & Operations $\to$ Representation $\to$ Cost & 为了缩小候选集提前维护了什么信息，更新/空间/I/O 付什么代价？\\
排序 & Operation + Invariant + Progress & 每一趟之后哪个区域已经确定，算法如何逐步扩大正确区域？\\
KMP & Representation of history & 怎样把已匹配历史压缩成 failure state，避免重复比较？\\
\bottomrule
\end{longtable}

\section{整门数据结构的专题地图}
\begin{longtable}{L{1.45cm}L{4.4cm}L{\dimexpr\linewidth-7.12cm\relax}}
\toprule
\textbf{编号} & \textbf{专题} & \textbf{它真正回答的问题}\\
\midrule
DS-01 & 复杂度与成本模型 & 输入规模怎样变成基本操作次数，怎样区分最坏/平均/摊还/I/O 成本？\\
DS-02 & 线性表与数组映射 & 连续位置与显式指针分别换来什么能力与代价？\\
DS-03 & 栈、队列与受控遍历 & 限制操作顺序怎样制造 LIFO/FIFO 与 frontier 语义？\\
DS-04 & 树与递归层次 & 层次关系怎样缩小候选空间并支持递归分治？\\
DS-05 & 堆与并查集 & 只维护目标所需的最少信息，怎样支持极值与动态集合？\\
DS-06 & 图表示与遍历 & 任意关系中怎样维护 Visited、Frontier 与 Expansion Policy？\\
DS-07 & 图算法 & MST、最短路、拓扑、关键路径分别维护什么“已确定”不变量？\\
DS-08 & 查找与有序索引 & 怎样用有序性、层次或索引减少候选集？\\
DS-09 & 散列 & 怎样把比较问题改造成 Key$\to$CandidateLocation 与冲突管理？\\
DS-10 & KMP & 怎样保存已匹配信息避免失败后完全重来？\\
DS-11 & 内部排序 & 怎样通过局部动作不断扩大已确定有序区域？\\
DS-12 & 外部排序 & 当 I/O 成为主成本时，怎样减少归并趟数与随机访问？\\
\bottomrule
\end{longtable}

\section{数据结构最重要的概念边界：不是背“$\neq$”，而是学判据}
\begin{longtable}{L{3.6cm}L{4.2cm}L{4.0cm}L{\dimexpr\linewidth-13.07cm\relax}}
\toprule
\textbf{概念边界} & \textbf{为什么容易混 / 真正区别} & \textbf{题目中的判定信号} & \textbf{混淆后的典型错误}\\
\midrule
\key{Logical Structure} $\neq$ \key{Storage Representation} & 一个描述“关系是什么”，一个描述“机器怎样保存关系” & 题目说“线性表”还不能推断是数组还是链表；必须看顺序/链式存储 & 把“线性表”直接等同于连续内存，插删/地址计算全错 \\
\midrule
\key{Structure} $\neq$ \key{Algorithm} & 结构是状态及 invariant，算法是改变状态的过程 & 问“AVL 是什么”与“插入后怎样旋转”分别属于对象定义和操作 & 背会代码但不知道何时需要修复、为什么修复 \\
\midrule
\key{Heap Order} $\neq$ \key{Total Order} & 堆只保证父子偏序，不保证兄弟/跨子树全局大小 & 题目若问“第二大在哪里/能否二分查找”要警惕 & 误以为堆数组整体有序 \\
\midrule
\key{Average $O(1)$} $\neq$ \key{Worst-case $O(1)$} & 平均复杂度依赖分布/装填等假设，最坏界是对所有输入保证 & Hash 题出现“平均”“最坏”“装填因子”“冲突” & 把散列表任何情况下都写成常数时间 \\
\midrule
\key{Shortest Path} $\neq$ \key{MST} & 一个优化“某些点之间的路径长度”，一个优化“连通全部顶点的总边权” & 问源点到各点距离看最短路；问连通所有点且总成本最小看 MST & 用 Prim/Kruskal 回答单源最短路 \\
\midrule
\key{Stable Sort} $\neq$ \key{In-place Sort} & 稳定性讨论相等 key 的相对次序；原地性讨论额外空间 & 题目同时给“相同关键字记录”和“辅助空间”时要分别判断 & 把“稳定”误认为“不额外开数组” \\
\midrule
\key{Asymptotic Cost} $\neq$ \key{Actual Machine Cost} & 渐近阶忽略常数和机器层次，真实时间还受 cache/I/O 等影响 & 外存、cache locality、数据规模很小时不能只看 $O(\cdot)$ & 认为同为 $O(n)$ 的数组/链表真实性能必然一样 \\
\bottomrule
\end{longtable}

\section{数据结构做题入口：这条箭头怎样实际使用}
\begin{examBox}{数据结构七步检查链}
\[
\boxed{\begin{gathered}
Goal\rightarrow Relation\rightarrow Operations/Workload\rightarrow Representation\\
\downarrow\\
Invariant\rightarrow State/Progress\rightarrow Cost/Boundary
\end{gathered}}
\]
\begin{enumerate}
  \item \key{Goal：}题目最后要求查找、更新、遍历、排序还是结构选择？
  \item \key{Relation：}对象是线性、层次、图、集合还是 key$\to$value 映射？
  \item \key{Operations/Workload：}哪些操作出现得最多，是否有最坏界、稳定性、范围查询等额外要求？
  \item \key{Representation：}题目实际用数组、指针、矩阵、邻接表、桶还是索引树？
  \item \key{Invariant：}当前结构合法必须满足什么？把它写在草稿边上。
  \item \key{State/Progress：}模拟算法时每一步哪些区域已经确定、frontier 在哪里、问题规模是否缩小？
  \item \key{Cost/Boundary：}最后数比较/移动/空间/I/O，并检查空结构、首尾、重复 key、不连通等边界。
\end{enumerate}
\end{examBox}

\begin{methodbox}{小例子：题目问“频繁插入并删除最大值，选什么结构？”}
不要看到“最大值”就直接答排序。按七步走：Goal=动态维护；Relation=元素集合本身没有必要完全有序；Operations=Insert/DeleteMax 高频；Representation 候选有有序数组、BST、Heap；Invariant 只需保证能快速找到极值；Cost 比较后，Heap 用较弱的父子偏序换取 $FindMax=O(1)$ 与更新 $O(\log n)$。这就是从需求推出结构，而不是从关键词猜结构。
\end{methodbox}

\section{一页压缩}
\begin{corebox}{闭眼后应该剩下什么}
\[
\boxed{\text{Problem}\rightarrow\text{Relation}\rightarrow\text{Operations}\rightarrow\text{Representation}\rightarrow\text{Invariant}\rightarrow\text{Cost}}
\]
快速查询几乎总来自提前维护更多信息；保存越多不是越好，只保存足以完成目标的关系；算法的本质是在局部修改后恢复/保持结构不变量。看到箭头时，把它读成“下一步我要问什么”，而不是背一串英文。
\end{corebox}

% ============================================================
\chapter{计算机组成原理：软件语义怎样成为硬件时间过程}
% ============================================================

\section{学科母问题}
\begin{corebox}{计组真正研究什么}
\key{如何用有限数量、有限速度、有限带宽的真实硬件，正确而高效地实现 ISA 对软件承诺的程序语义？}
\end{corebox}

一条指令从软件看可以抽象为状态转移 $S_{t+1}=F_I(S_t)$。这里 $S_t$ 指程序员可见的体系结构状态，例如 PC、通用寄存器和可见内存；$F_I$ 表示当前指令 $I$ 规定的语义。计组真正要解释的是：\key{这个看起来“一步完成”的语义，怎样在硬件内部被拆成若干数据移动和时间阶段。}

母模型：
\[
\boxed{\text{ISA State}\rightarrow\text{Data Location}\rightarrow\text{Datapath/Control}\rightarrow\text{Timing}\rightarrow\text{Commit}\rightarrow\text{Cost}}
\]

\subsection{逐格解释}
\begin{description}
  \item[ISA State] 指令执行前后，软件能够观察到哪些状态变化；这是硬件必须兑现的承诺。
  \item[Data Location] 本次运算所需指令和操作数当前在寄存器、Cache、主存还是设备中。
  \item[Datapath / Control] 数据实际经过 ALU、MUX、寄存器、总线、Memory interface 的路径；控制信号决定本周期哪些路径被启用。
  \item[Timing] 每个值在哪个周期产生、什么时候稳定、什么时候能被后续部件使用。
  \item[Commit] 内部可能预测、并行、暂存，但最终什么时候才能把结果作为“这条指令已经完成”的体系结构结果暴露。
  \item[Cost] 关键路径、cycle、CPI、stall、miss penalty、带宽等性能代价。
\end{description}
箭头表示\key{分析顺序}：先知道指令应该做什么，才能判断数据该去哪；只有路径明确以后，才能正确谈时序和性能。

\section{贯穿例子：一条 LOAD 怎样映射到整门计组}
以 \texttt{LOAD R1, 8(R2)} 为例。这里不展开具体 ISA 编码，只用它导航整门学科。

\begin{longtable}{L{3.0cm}L{\dimexpr\linewidth-3.85cm\relax}}
\toprule
\textbf{母模型位置} & \textbf{LOAD 中实际发生什么}\\
\midrule
ISA State & 语义目标：$R1\leftarrow M[R2+8]$，并按 ISA 规则推进 PC。\\
Data Location & 指令要先被取到；$R2$ 在寄存器；目标数据可能在 L1/L2/DRAM。\\
Datapath / Control & 读寄存器 $R2$，ALU 计算有效地址，选择数据存储端口，返回数据后选择写回 $R1$。\\
Timing & 地址在哪一级产生？Cache hit 数据何时可用？若后续指令马上读 $R1$，需要 forwarding 还是 stall？\\
Commit & 只有符合微架构/ISA 提交规则后，$R1$ 的新值才成为软件可见结果。\\
Cost & Cache miss、结构冲突、数据依赖都会增加等待；性能不能只由时钟频率决定。\\
\bottomrule
\end{longtable}

于是：指令系统讲“LOAD 承诺什么”；数据通路讲“怎么走”；流水线讲“什么时候走”；Cache 讲“数据在哪里”；性能讲“等了多久”。这就是总图中“一个母例题串全科”的意义。

\section{四个坐标轴，以及它们对应的章节}
\subsection{Semantic State：什么对软件可见？}
对应：数据表示、ISA、指令系统、异常/中断的体系结构语义。学习时问：同一串 bit 在当前语义下代表什么？这条指令最终允许改变哪些可见状态？

\subsection{Data Location：需要的数据在哪里？}
对应：寄存器、主存、Cache、I/O。学习时问：目标数据当前在哪一级？如果不在，下一层是谁？搬运成本是什么？

\subsection{Datapath / Control / Timing：硬件如何在时间上组织工作？}
对应：单周期、多周期、流水线、硬布线/微程序控制。Datapath 决定有哪些路，Control 决定此刻打开哪条路，Clock 决定状态何时采样。

\subsection{Resource / Cost：哪里会等待？}
对应：hazard、总线仲裁、Cache miss、I/O、性能计算。ALU、端口、总线和存储系统都有限；“值没准备好”和“资源被占用”都会变成 stall。

\section{母模型—章节映射}
\begin{longtable}{L{3.0cm}L{4.2cm}L{\dimexpr\linewidth-8.47cm\relax}}
\toprule
\textbf{章节/专题} & \textbf{主要对应母模型位置} & \textbf{学习时真正要问}\\
\midrule
数据表示与运算 & ISA State / Semantics & 有限 bit 如何解释，溢出/舍入条件是什么？\\
指令系统/寻址 & ISA State $\to$ Data Location & 指令承诺什么，操作数如何被定位？\\
主存组织 & Data Location / Resource & 地址线、数据线、bank 怎样决定容量与并行访问？\\
数据通路与控制 & Datapath / Control & 一条指令被拆成哪些微操作，哪些控制信号驱动？\\
流水线 & Timing / Availability & 每个值在哪一级产生，hazard 为什么迫使等待/转发/预测？\\
Cache & Data Location / Fast-Slow Path & 目标块放哪、如何识别、miss 后去哪、写策略怎样维护副本？\\
总线与 I/O & Resource / Interface & 多部件如何共享通信资源，异步设备怎样通知/搬运？\\
性能 & Cost / Bottleneck & IC、CPI、周期时间、miss/stall 怎样组合成完成时间？\\
\bottomrule
\end{longtable}

\section{计组专题地图}
\begin{longtable}{L{1.45cm}L{4.45cm}L{\dimexpr\linewidth-7.17cm\relax}}
\toprule
\textbf{编号} & \textbf{专题} & \textbf{核心问题}\\
\midrule
CO-01 & 数据表示与运算 & 有限位宽怎样表达整数、浮点与运算语义，溢出/舍入从哪里来？\\
CO-02 & ISA 与寻址方式 & 软件可见机器承诺是什么，一条指令怎样描述操作与操作数？\\
CO-03 & 主存与芯片组织 & 容量、位宽、bank 和地址线怎样落成实际存储系统？\\
CO-04 & 数据通路与控制器 & 指令如何被分解成微操作，控制信号如何驱动数据移动？\\
CO-05 & 单周期、多周期与流水线 & 工作怎样沿时间切分与重叠，hazard 为什么出现？\\
CO-06 & Cache 与存储层次 & 如何利用局部性管理副本，placement/tag/replacement/write 各解决什么？\\
CO-07 & 总线与 I/O & 同步 CPU 如何与异步设备通信，Polling/Interrupt/DMA 分别卸载什么？\\
CO-08 & 性能与并行 & IC、CPI、周期时间、瓶颈、并行度怎样共同决定完成时间？\\
\bottomrule
\end{longtable}

\section{计组最重要的概念边界：用题目信号判断}
\begin{longtable}{L{3.6cm}L{4.3cm}L{4.0cm}L{\dimexpr\linewidth-13.17cm\relax}}
\toprule
\textbf{概念边界} & \textbf{真正区别} & \textbf{题目信号} & \textbf{混淆后的错误}\\
\midrule
\key{ISA} $\neq$ \key{Microarchitecture} & ISA 是软件可见承诺；微架构是内部实现方法 & 问“指令/寄存器/寻址语义”看 ISA；问“流水级/预测/内部端口”看实现 & 把某 CPU 的实现细节当成所有同 ISA 机器都必须如此 \\
\midrule
\key{CPI 低} $\neq$ \key{程序必然快} & 总时间还由 IC 和 cycle time 决定 & 比较机器时题目常同时给 IC/CPI/频率 & 只比较 CPI 得出性能结论 \\
\midrule
\key{Pipeline Throughput} $\neq$ \key{Instruction Latency} & 流水线主要提高单位时间完成指令数，不代表单条指令经过的级数减少 & 问“填满后每周期完成几条”与“某条从 IF 到 WB 多久”是不同量 & 把五级流水说成每条指令只需 1 cycle \\
\midrule
\key{Address Field} $\neq$ \key{Full Address} & 指令字段可能是寄存器号、立即数、偏移，需要寻址方式计算有效地址 & 看到 base+offset、间接寻址、PC-relative & 直接把指令中的字段当绝对地址 \\
\midrule
\key{Cache Miss} $\neq$ \key{Page Fault} & 一个是硬件存储层次副本未命中；一个是虚拟内存映射/驻留异常，进入 OS & 题目分别给 tag/index 与 PTE/present 位 & Cache miss 后错误地直接认为要访问磁盘/陷入 OS \\
\midrule
\key{Interrupt} $\neq$ \key{DMA} & Interrupt 解决“完成后怎样通知 CPU”；DMA 解决“数据怎样批量搬运” & 题目问 CPU 是否逐字节搬、谁产生完成通知 & 把 DMA 理解成不需要中断/CPU完全不参与 \\
\midrule
\key{Clock Rate 高} $\neq$ \key{Workload 更快} & 高频率可能伴随更高 CPI、更多指令或不同 miss 行为 & 真正比较同一程序完成时间时必须综合指标 & 只看 GHz 判断处理器快慢 \\
\bottomrule
\end{longtable}

\section{计组做题入口：英文链怎样变成草稿动作}
\begin{examBox}{计组六步检查链}
\[
\boxed{State\rightarrow DataLocation\rightarrow Path\rightarrow Resource\rightarrow Timing\rightarrow Commit/Cost}
\]
\begin{enumerate}
  \item \key{State：}题目执行前 PC、寄存器、Cache/Memory 关键状态是什么？目标指令最终应该改变什么？
  \item \key{DataLocation：}指令和操作数当前在哪一级；若 miss，下一层在哪里？
  \item \key{Path：}画最小数据通路，只标当前题真正经过的部件和 MUX。
  \item \key{Resource：}同一周期是否有人争 ALU、端口、总线或存储器？值是否还没准备好？
  \item \key{Timing：}逐 cycle 标值何时产生、何时可用，必要时画流水表。
  \item \key{Commit/Cost：}确认最终可见状态，然后再算 CPI、时间、命中/未命中代价。
\end{enumerate}
\end{examBox}

\begin{methodbox}{小例子：LOAD 后紧跟 ADD 使用 R1}
先别背“数据冒险要停几拍”。State：LOAD 要写 R1；DataLocation：数据要到 MEM/Cache 后才产生；Path：LOAD 经过 EX 算地址、MEM 取数据，ADD 在自己的 EX 级需要 R1；Resource/Timing：比较“数据产生时刻”和“ADD 需要时刻”；若 forwarding 仍来不及才 stall。这样停几拍是从时间关系推出来的，而不是死背固定数字。
\end{methodbox}

\section{一页压缩}
\begin{corebox}{闭眼后应该剩下什么}
\[
\boxed{\text{ISA State}\rightarrow\text{Data Location}\rightarrow\text{Datapath/Control}\rightarrow\text{Timing}\rightarrow\text{Commit}\rightarrow\text{Cost}}
\]
硬件可以内部重叠、缓存、预测和复用，但不能改变 ISA 语义；性能问题先追踪数据位置、可用时刻和等待来源，再谈 CPI 或频率。
\end{corebox}

% ============================================================
\chapter{操作系统：有限硬件怎样成为受保护的共享世界}
% ============================================================

\section{学科母问题}
\begin{corebox}{OS 真正研究什么}
\key{当多个彼此独立甚至互不信任的程序共享 CPU、内存、设备和持久存储时，内核怎样建立抽象、保持控制、协调等待与竞争、隔离权限，并在事件和故障中维持合法状态？}
\end{corebox}

OS 的最简状态模型：
\[
\boxed{S=(Objects,Relations,Queues)}.
\]
事件到来以后：
\[
S_t\xrightarrow{Event+Mechanism+Policy}S_{t+1},
\]
同时要求 Invariant 成立，并优化 Cost。

\subsection{这些词分别是什么意思}
\begin{description}
  \item[Objects] 内核正在管理的“东西”，例如 Process、Thread、AddressSpace、Page、File、OpenFile、DeviceRequest。
  \item[Relations] 对象之间怎样互相引用或映射，例如 fd$\to$OpenFile$\to$inode，VPN$\to$PTE$\to$PFN，父进程$\to$子进程。
  \item[Queues] 当前不能继续的请求/实体在哪里等待，例如 ready queue、I/O wait queue、lock waiters、device request queue。
  \item[Event] 真正触发变化的事情，例如 timer interrupt、page fault、I/O completion、unlock、system call。
  \item[Mechanism] 内核/硬件“有能力怎样做”，例如 context switch、trap entry、wakeup、page-table update。
  \item[Policy] 当存在多个合法选择时“应该选谁/什么时候做”，例如 scheduler 选谁、replacement 牺牲哪页、I/O scheduler 排哪个请求。
  \item[Invariant] 任何合法状态都必须满足的规则，例如权限不能越界、互斥不能同时被两者占用、页表映射必须符合保护位。
  \item[Cost] 等待、切换、TLB/cache miss、I/O、竞争、写回等代价。
\end{description}

箭头 $S_t\xrightarrow{E}S_{t+1}$ 的最重要含义是：\key{没有事件，不要凭感觉改状态；有事件，还要说清是谁用什么机制改了哪一部分状态。}

\section{贯穿例子：一次阻塞 read 怎样把模型跑起来}
假设进程 A 调用 \texttt{read()}，目标数据不在页缓存，需要磁盘读取。

\begin{longtable}{L{3.0cm}L{\dimexpr\linewidth-3.85cm\relax}}
\toprule
\textbf{模型位置} & \textbf{这个例子里具体是什么}\\
\midrule
Objects & Process A、fd/OpenFile/inode、page-cache page、I/O request、device。\\
Relations & A 的 fd 指向 OpenFile，OpenFile 指向文件对象；文件偏移映射到目标数据页/块。\\
Queues & A 因等待数据进入 wait queue；磁盘请求进入设备请求队列；其他进程仍可能在 ready queue。\\
Event & \texttt{read()} 系统调用、cache miss、I/O completion interrupt。\\
Mechanism & trap 进入内核、提交 I/O、把 A 置 Blocked、completion 时 wakeup 成 Ready。\\
Policy & A 阻塞后 scheduler 选哪个 Ready 进程；A 被唤醒后何时再次得到 CPU。\\
Invariant & Blocked 进程不能被当成 Ready；权限/文件引用必须合法；I/O 数据完成后才能向用户返回。\\
Cost & 系统调用、context switch、设备 I/O 和等待时间远高于纯内存 fast path。\\
\bottomrule
\end{longtable}

这个例子分别连接进程/调度、文件系统、Page Cache、I/O、中断；总图只负责告诉你\key{这些专题在同一条状态轨迹里在哪里交接}，具体细节由专题册拥有。

\section{OS 的五个不可替代职责，以及它们出现在哪里}
\begin{longtable}{L{2.8cm}L{5.0cm}L{\dimexpr\linewidth-9.07cm\relax}}
\toprule
\textbf{职责} & \textbf{含义} & \textbf{典型章节/例子}\\
\midrule
Control & 用户代码可以直接高效执行，但内核必须能受控进入/收回 CPU & user/kernel mode、system call、trap、timer interrupt \\
Virtualization & 把有限复杂硬件包装成稳定私有/逻辑抽象 & process/thread、virtual address space、file/fd \\
Coordination & 多个执行流争 CPU、锁、buffer、device 时需要排队与唤醒 & scheduling、PV/lock、I/O wait、device queue \\
Protection & 规定“谁能对哪个对象做什么”并让硬件可执行地检查 & privilege、page permissions、file permissions、isolation \\
Persistence & 状态跨运行时存在，多步更新失败后仍要有可恢复语义 & file system、writeback、fsync、journaling/crash consistency \\
\bottomrule
\end{longtable}

\section{OS 专题地图}
\begin{longtable}{L{1.45cm}L{4.45cm}L{\dimexpr\linewidth-7.17cm\relax}}
\toprule
\textbf{编号} & \textbf{专题} & \textbf{核心问题}\\
\midrule
OS-01 & 控制权、系统调用与中断异常 & 用户代码怎样安全进入/离开内核，mode switch 与 context switch 如何区分？\\
OS-02 & 进程、线程与调度 & 执行实体如何创建、阻塞、唤醒、被调度并最终退出？\\
OS-03 & 并发、锁与 PV & 共享状态面对任意交错时怎样保持 safety、ordering 与 progress？\\
OS-04 & 虚拟内存与地址翻译 & VA 怎样映射到 PA，fault、TLB、COW、replacement 各解决什么？\\
OS-05 & I/O 系统 & 快 CPU 怎样接入慢而异步的设备，Interrupt、DMA、Buffer/Cache 各自作用是什么？\\
OS-06 & 文件系统 & Name、fd、open state、inode 与 block 怎样解耦并形成可命名持久对象？\\
OS-07 & 持久化与崩溃一致性 & 多步写入中途失败时怎样保证或恢复合法持久状态？\\
OS-B1 & CPU $\times$ OS 桥梁 & 特权级、MMU、TLB、中断、Timer、DMA 的软硬件分工边界是什么？\\
OS-I1 & OS 综合状态机 & read/mmap/fork/COW/pipe/memory pressure 等真实过程怎样跨多个专题运行？\\
\bottomrule
\end{longtable}

\section{母模型—专题映射：每一章在改 S 的哪一部分}
\begin{longtable}{L{3.0cm}L{4.25cm}L{\dimexpr\linewidth-8.52cm\relax}}
\toprule
\textbf{专题} & \textbf{主要改变/维护的状态} & \textbf{用母模型观察时问什么}\\
\midrule
进程/调度 & 进程对象 + Ready/Blocked 等队列 & 什么事件改变资格？wakeup 与 schedule 谁负责？\\
并发/PV & 共享对象 + 锁/信号量状态 + 等待队列 & 哪个 invariant 被交错威胁？谁持有、谁等待、谁能制造条件？\\
虚拟内存 & VA$\to$PTE$\to$frame relations + residency state & 这次访问是 translation、permission 还是 residency 问题？\\
I/O & I/O 请求 + 设备队列 + 完成事件 & 请求提交后谁等？数据谁搬？完成事件怎样把进程变 Ready？\\
文件系统 & Name、fd、OpenFile、inode 与 block 的引用/映射关系 & 名字、打开实例、文件实体、数据位置分别是哪一层对象？\\
持久化 & 内存状态 $\leftrightarrow$ 稳定介质状态 & write 返回到哪一步？crash point 后哪些状态允许存在？\\
\bottomrule
\end{longtable}

\section{OS 需要牢牢记住的四条横向规律：每条都要会落地}
\subsection{对象不是名字，引用与映射决定关系}
PID、fd、虚拟地址、路径名都不是对象本身。例如两个 fd 可以指向同一个 OpenFile，也可以通过两个独立 OpenFile 最终指向同一 inode。题目一旦问“共享 offset 吗”，就不能只看文件名相同，而要追踪引用链。

\subsection{Blocked 与 Ready 之间隔着“事件完成”，Ready 与 Running 之间隔着“调度”}
磁盘完成会把等待进程从 Blocked 变 Ready；它只获得竞争 CPU 的资格。只有 scheduler 选中后才 Ready$\to$Running。做 I/O/调度综合题时，这两根箭头必须分开画。

\subsection{Mechanism 与 Policy 严格分开}
Context switch 是“能切换”的机制，scheduler 选谁是策略；page fault handler 能回收/建立映射是机制，replacement 选 victim 是策略。题目问“由谁决定下一个进程”时，不要回答“上下文切换”。

\subsection{Fast Path 失败时才进入 Kernel Slow Path}
TLB hit、page-cache hit、无竞争锁等常见情况尽量便宜；fault、I/O、锁竞争、回收与恢复进入慢路径。性能题要先判断“这次到底命中了哪条路径”，再累加代价。

\section{OS 最重要的概念边界：判断题目到底在问哪一个状态}
\begin{longtable}{L{3.6cm}L{4.35cm}L{3.95cm}L{\dimexpr\linewidth-13.17cm\relax}}
\toprule
\textbf{概念边界} & \textbf{真正区别} & \textbf{题目信号} & \textbf{混淆后的错误}\\
\midrule
\key{Process} $\neq$ \key{Program} & Program 是静态代码/数据；Process 是运行中的执行与资源状态 & 看到 PID、状态、地址空间、调度就是 process & 认为同一程序只能有一个进程 \\
\midrule
\key{Ready} $\neq$ \key{Blocked} & Ready 只缺 CPU；Blocked 还缺某个外部条件/资源 & “时间片到”通常 Running$\to$Ready；“等待 I/O/信号量”是 Blocked & I/O 完成后直接写 Running，跳过调度 \\
\midrule
\key{Mode Switch} $\neq$ \key{Context Switch} & mode switch 可仍是同一执行实体；context switch 改变 CPU 当前执行实体 & system call/interrupt 不一定换进程；schedule 才涉及执行实体切换 & 认为每次系统调用都一定切换进程 \\
\midrule
\key{Mutex} $\neq$ \key{Condition} & Mutex 保护谁能同时进入；Condition 表达“什么时候才允许继续” & 看到共享临界区看互斥；看到“必须等到 buffer 非空”看条件同步 & 用一把锁代替状态条件，或持锁等待造成死锁 \\
\midrule
\key{Signal/Wakeup} $\neq$ \key{Predicate True} & 通知只是“状态可能变化”；真正条件仍要重新检查 & condition variable 常用 while(predicate false) wait & 被唤醒后不检查条件直接继续 \\
\midrule
\key{TLB Miss} / \key{Page Fault} / \key{Cache Miss} & 翻译缓存、虚拟页映射/驻留、数据缓存是三个不同层次状态 & 看 TLB / PTE / present / cache tag 各自给了什么 & 一看到 miss 就默认磁盘 I/O \\
\midrule
\key{Buffer} $\neq$ \key{Cache} $\neq$ \key{Spooling} & Buffer 吸收速率/粒度差；Cache 为复用保存副本；Spooling 用队列虚拟化独占设备 & 题目问“平滑生产消费”“重复访问”“打印排队”分别不同 & 把任何内核内存都叫缓存，机制目的判断错 \\
\midrule
\key{File Name} $\neq$ \key{inode} $\neq$ \key{Open File State} & Name 属命名空间；inode 是文件实体元数据；OpenFile 是一次打开的运行时状态（如 offset） & hard link/unlink 看 name$\leftrightarrow$inode；dup/fork offset 看 OpenFile & 认为删除名字就立即销毁对象，或两次 open 必共享 offset \\
\midrule
\key{write() return} / \key{Durable} & 返回通常表示数据已被内核接受，不等于稳定介质已经持久化 & 题目出现 page cache / dirty / fsync / crash & 把系统调用返回时刻当作物理落盘时刻 \\
\bottomrule
\end{longtable}

\section{OS 做题入口：九个英文词怎样变成实际推演}
\begin{examBox}{OS 九步检查链}
\[
\boxed{\begin{gathered}
Resource/Abstraction\rightarrow Objects/Relations/Queues\rightarrow Event\rightarrow Boundary\\
\downarrow\\
Mechanism\rightarrow Policy\rightarrow Invariant\rightarrow Wait/Fault/Recovery\rightarrow Cost
\end{gathered}}
\]
\begin{enumerate}
  \item \key{Resource/Abstraction：}底层真实资源是什么，上层看到的抽象是什么？
  \item \key{Objects/Relations/Queues：}画出题目真正涉及的对象、引用/映射和等待队列。
  \item \key{Event：}逐步写“谁发生了什么”，不要无事件改变状态。
  \item \key{Boundary：}当前是 user/kernel、hardware/kernel、name/object、Ready/Blocked 哪个边界？
  \item \key{Mechanism：}系统具备用什么机制完成变化？trap、switch、wakeup、PTE update、DMA？
  \item \key{Policy：}如果有多个候选，策略选谁？没有选择问题就不要硬找 policy。
  \item \key{Invariant：}权限、互斥、容量、引用生命周期、映射合法性有哪些必须保持？
  \item \key{Wait/Fault/Recovery：}正常路径失败后，是等待、fault、回收还是恢复？
  \item \key{Cost：}最后计算等待、切换、访存/I/O 次数等代价。
\end{enumerate}
\end{examBox}

\begin{methodbox}{小例子：I/O 完成后进程是什么状态？}
题目说 A 因读磁盘阻塞，之后磁盘中断到达。Objects/Queues：A 在 I/O wait queue；Event=I/O completion；Mechanism=wakeup，把 A 移入 ready queue；Policy 只有在调度时才决定是否选 A；Invariant=Blocked/Ready 语义不能混。因此答案首先是 Ready，而不是凭“数据到了”直接写 Running。这个判断正是九步链在一道选择题里的缩写版本。
\end{methodbox}

\section{一页压缩}
\begin{corebox}{闭眼后应该剩下什么}
\[
\boxed{Objects+Relations+Queues\xrightarrow{Event+Mechanism+Policy}NewState}
\]
做题时把这句话翻译成：\key{谁是什么对象？它和谁什么关系？谁在等？发生了什么？系统靠什么改变状态？如果有多个合法选择由什么策略选？改变后是否仍合法？} 这比背九个英文词更重要。
\end{corebox}

% ============================================================
\chapter{计算机网络：没有全局即时状态时怎样协作}
% ============================================================

\section{学科母问题}
\begin{corebox}{网络真正研究什么}
\key{大量自治节点只能看到局部状态，信息传播有时延、链路可能不可靠、资源还要共享时，怎样仅靠消息交换共同形成可用的端到端通信？}
\end{corebox}

网络统一模型：
\[
\boxed{\text{Scope}\rightarrow\text{Local State}\xrightarrow{\text{Message/Timer}}\text{Transition}\rightarrow\text{Feedback}\rightarrow\text{Cost}}.
\]

\subsection{逐格解释}
\begin{description}
  \item[Scope] 当前决策负责多远：一个应用语义、两个端点、跨多个网络、当前一跳，还是一段物理介质？
  \item[Local State] 做决定的设备现在真正知道什么；网络中没有一个节点可以瞬时读取所有其他节点状态。
  \item[Message / Timer] 网络状态机最常见两类输入：收到消息，或者“等了足够久仍没等到消息”。
  \item[Transition] 设备依据本地状态和输入更新表项、窗口、序号、队列或连接状态，并决定发送/丢弃/重传等动作。
  \item[Feedback] 本节点的动作通过新报文影响别的节点，使分布式系统逐步协调。
  \item[Cost] 时延、带宽、队列、丢包、在途数据量、重传和控制开销。
\end{description}
箭头表示\key{分析一个网络机制时的观察顺序}，不是说真实 packet 一定按这六个词逐层“走”。例如 Scope 是分析者先确定的坐标，不是网络中真的存在一个叫 Scope 的协议层。

\section{贯穿例子：TCP 超时重传怎样落进母模型}
\begin{longtable}{L{2.8cm}L{\dimexpr\linewidth-3.65cm\relax}}
\toprule
\textbf{母模型位置} & \textbf{TCP 超时场景}\\
\midrule
Scope & 端到端：发送方和接收方之间的一条 TCP 连接。\\
Local State & 发送方保存已发送未确认的序号范围、重传计时器、窗口等；路由器通常不保存这些 TCP 序号状态。\\
Message/Timer & 正常时 ACK 到达；异常时 timer expiration 表示发送方在预期时间内没有获得足够反馈。\\
Transition & 发送方更新重传/拥塞相关状态并重新发送相应 segment。\\
Feedback & 重传数据到达接收方，接收方 ACK 又返回发送方，继续改变其状态。\\
Cost & 超时增加至少一个等待周期，并消耗额外带宽；窗口/拥塞策略还会影响后续吞吐。\\
\bottomrule
\end{longtable}
这就是“网络是耦合的局部状态机”最小实例。学习 GBN、SR、TCP、路由收敛和拥塞时，都应该尝试找出各自的 Local State、Message/Timer 和 Feedback。

\section{六个现实约束怎样生成网络机制}
\begin{longtable}{L{2.8cm}L{3.8cm}L{4.1cm}L{\dimexpr\linewidth-11.97cm\relax}}
\toprule
\textbf{约束} & \textbf{直接困难} & \textbf{机制方向} & \textbf{在哪些章节看到}\\
\midrule
Distance & 远端状态不能瞬间知道，传播要时间 & RTT、Timer、Pipeline、Window & 性能、可靠传输、TCP \\
Unreliability & bit/packet 可能错、丢、重、乱 & CRC/Checksum、Seq、ACK、Retransmission & 链路差错、可靠传输 \\
Sharing & 多个流共享同一介质/链路/队列 & MAC、Multiplexing、Queue、Congestion & LAN/MAC、拥塞 \\
Heterogeneity & 不同介质、设备与局域网技术不统一 & Layering、统一网络层抽象、IP & 体系结构、IP \\
Scale & 全球节点太多，不能保存逐主机平坦状态 & Prefix、Hierarchy、Aggregation、AS & CIDR、路由、BGP \\
No Global Instantaneous State & 每个节点只能用延迟到达的局部信息推断 & Distributed Routing、Timeout、Feedback、Convergence & 路由、TCP、拥塞 \\
\bottomrule
\end{longtable}

这张表的使用方式是：学到一个新机制时倒着问“如果没有左边那个现实约束，这个机制还是否必要？”例如没有距离和 RTT，Stop-and-Wait 的低利用率问题就会弱很多；没有共享资源，拥塞控制就失去核心动机。

\section{第一坐标：Scope——一个决定负责多远}
\begin{longtable}{L{2.5cm}L{4.4cm}L{\dimexpr\linewidth-8.17cm\relax}}
\toprule
\textbf{作用域} & \textbf{核心问题} & \textbf{典型对象/状态与题目例子}\\
\midrule
Application & 双方交换的语义是什么？ & HTTP method/status、DNS query/answer；问“请求含义/应用交互”在这里 \\
Transport & 到哪个端点、需要什么端到端通信条件？ & Port、connection、Seq/ACK、window；问“进程、可靠、流控”在这里 \\
Network & 跨网络目的地在哪里，当前节点下一跳去哪？ & IP、prefix、forwarding state；问“子网、LPM、路由”在这里 \\
Link & 当前这一跳怎样交给下一设备？ & Frame、MAC、CRC、switch state；问“交换机、CSMA、帧”在这里 \\
Physical & bit 怎样成为介质中的信号？ & Symbol、encoding、modulation、bandwidth；问 Nyquist/Shannon/码元在这里 \\
\bottomrule
\end{longtable}

纵向通信主链：
\[
\boxed{\text{Meaning}\rightarrow\text{Endpoint}\rightarrow\text{Path}\rightarrow\text{One Hop}\rightarrow\text{Signal}}.
\]
这条箭头的意思是：\key{一个应用消息要真正传出去，必须逐步补齐“端点是谁、跨网目的地是谁、这一跳交给谁、最后怎样编码成信号”的信息。} 封装就是为不同作用域加入完成局部决策所需的控制信息，而不是为了背“每层加首部”。

\section{第二坐标：State Ownership——谁保存完成决策所需的信息}
\begin{itemize}
  \item TCP connection/Seq/ACK/window：主要属于通信端点；
  \item switch：维护 MAC$\to$Port；
  \item router data plane：维护 Prefix$\to$NextHop/Interface；
  \item routing protocol：维护形成/更新上述转发表所需的控制状态；
  \item DNS resolver/server：维护 Name$\to$ResourceRecord 与 TTL/authority 相关状态。
\end{itemize}

\begin{methodbox}{“谁保存状态”怎样用于题目}
如果题目问“某 TCP ACK 到达后谁更新发送窗口”，先问状态所有者：TCP endpoint。普通中间路由器只需处理 IP forwarding，不会因为看到了 TCP ACK 就维护该连接的 seq / window。反过来，题目问“目的 IP 最长前缀匹配哪一项”，状态所有者是当前路由器的 forwarding table，不要跑去分析发送方 TCP 状态。
\end{methodbox}

\section{三条贯穿网络的关系链：每条都给出使用场景}
\subsection{Name $\rightarrow$ Address $\rightarrow$ Forwarding $\rightarrow$ Delivery}
\[
\boxed{\text{Name}\rightarrow\text{IP Address}\rightarrow\text{Forwarding Decision}\rightarrow\text{NextHop}\rightarrow\text{Link Address}}.
\]
中文意思：用户/应用先用名字描述“我要找谁”；DNS 等机制得到网络层地址；当前节点用目的 IP 和转发表决定下一跳；进入具体链路后，还要得到下一跳的链路层地址。做“浏览器访问网站”“ARP 与路由器”“同网段/异网段”综合题时，用这条链逐层追踪。

\subsection{Reliability = Sequence + ACK + Timer + Retransmission + Window}
“$+$”表示可靠传输需要这些信息/机制共同配合。Sequence 用于区分新旧/顺序，ACK 提供成功反馈，Timer 处理“没有反馈”的情况，Retransmission 修复丢失，Window 允许多个数据在途提高利用率。链路层 GBN/SR 与 TCP 都会调用这套母模型，但作用域和具体规则不同。

\subsection{Statistical Multiplexing $\rightarrow$ Queue $\rightarrow$ Congestion $\rightarrow$ Feedback}
多个流共享资源提高平均利用率；突发到达超过服务能力时，请求先排队；持续过载使延迟/丢包上升形成拥塞；端点只能通过丢包、RTT、显式信号等反馈调节负载。做拥塞题时不要从 cwnd 公式开始，而要先找“哪个有限资源正在排队”。

\section{必须三分：Reliability / Flow / Congestion}
\begin{longtable}{L{2.8cm}L{2.7cm}L{3.6cm}L{\dimexpr\linewidth-10.37cm\relax}}
\toprule
\textbf{机制} & \textbf{保护对象} & \textbf{做题时问什么} & \textbf{典型信号}\\
\midrule
Reliability & 通信语义 & 数据有没有正确、完整地到达？ & seq、ACK、timeout、duplicate、retransmit \\
Flow Control & Receiver & 接收端缓冲和处理能力能不能承受？ & receive window、receiver buffer、rwnd \\
Congestion Control & Network & 中间路径共享资源能不能承受？ & cwnd、loss/RTT、slow start、AIMD、queue \\
\bottomrule
\end{longtable}
三者都可能让“发送量变小”，所以容易混；真正区分标准不是动作，而是\key{它在保护谁的状态}。

\section{Forwarding 与 Routing 必须彻底分开}
\begin{methodbox}{Data Plane / Control Plane：一个用表，一个造表}
\textbf{Forwarding：}当前 packet 已到某台路由器，使用现有 forwarding table 做局部快速决策。
\[
DestinationIP\xrightarrow{LPM}NextHop.
\]
这里箭头表示“输入目的地址，通过最长前缀匹配查表，输出下一跳”。

\textbf{Routing：}多个路由器怎样交换控制信息，让 forwarding table 被建立或更新。
\[
LocalKnowledge+RoutingMessages\rightarrow UpdatedRoutingState.
\]
题目若给一张路由表让你查下一跳，是 forwarding；题目若让你迭代 DV、运行 Dijkstra、分析 RIP/OSPF/BGP，是 routing/control plane。
\end{methodbox}

\section{母模型—章节映射}
\begin{longtable}{L{3.0cm}L{4.25cm}L{\dimexpr\linewidth-8.52cm\relax}}
\toprule
\textbf{章节/专题} & \textbf{主要对应母模型位置} & \textbf{学习时真正要问}\\
\midrule
物理/性能 & Cost + Distance & 一 bit/packet 的发送、传播、排队和吞吐代价怎么形成？\\
链路/MAC & One-hop Scope + Local State & 同一介质上谁能发、帧怎样一跳交付、错误怎样检测？\\
可靠传输 & Message/Timer + Feedback & 数据/ACK 丢失或延迟时，端点状态机怎样修复？\\
IP/转发 & Network Scope + Forwarding State & 目的地址如何被 prefix match 映射成 next hop？\\
路由 & Distributed Local State + Feedback & 每个路由器只靠局部消息怎样逐步形成可用路径知识？\\
TCP/UDP & Endpoint Scope + Connection State & 主机间 IP 服务怎样变成进程间通信与连接状态？\\
拥塞 & Queue + Feedback + Cost & 需求超过路径能力时，端点怎样从有限反馈调整负载？\\
DNS/HTTP 等应用 & Meaning / Name & 名字怎样变地址，bytes 怎样获得应用语义？\\
\bottomrule
\end{longtable}

\section{网络专题地图}
\begin{longtable}{L{1.45cm}L{4.45cm}L{\dimexpr\linewidth-7.17cm\relax}}
\toprule
\textbf{编号} & \textbf{专题} & \textbf{唯一母问题}\\
\midrule
NW-01 & 通信基础与网络性能 & 一个 bit 穿过真实信道要付什么时间和容量代价？\\
NW-02 & 单跳交付：帧、MAC、LAN & 下一设备已知后，这一跳怎样共享介质、检测错误并完成交付？\\
NW-03 & 可靠传输 & 怎样用 Seq/ACK/Timer/Retransmission/Window 在不可靠信道上构造可靠服务？\\
NW-04 & IP 地址、子网与分组转发 & 目的地址怎样通过 Prefix Match 变成下一跳？\\
NW-05 & 路由与控制平面 & 转发表从哪里来，局部知识怎样通过消息逐步收敛？\\
NW-06 & 运输层与 TCP 状态机 & IP 到主机后怎样交给进程，并建立端点之间的连接状态？\\
NW-07 & 拥塞与反馈控制 & Demand$>$Capacity 时怎样通过反馈调节在途负载？\\
NW-08 & 应用层：发现与语义 & 底层能传 bytes 以后，怎样解决 Name$\to$Address 与 Bytes$\to$Meaning？\\
NW-I1 & 一个网络请求的一生 & DHCP/DNS/ARP/IP/TCP/HTTP 等机制在真实访问中怎样接起来？\\
NW-B1 & OS $\times$ Network 桥梁 & socket、协议栈、buffer、NIC、DMA、中断和进程唤醒怎样交接？\\
\bottomrule
\end{longtable}

\section{一个网络请求的一生：总图中的箭头怎样读}
第一次访问远程网站时，可用下面链条作为\key{专题导航器}：
\[
\boxed{\begin{gathered}
NetworkIdentity\rightarrow NameResolution\rightarrow NextHopDecision\rightarrow OneHopDelivery\\
\downarrow\\
Routing/Forwarding\rightarrow TransportEndpoint\rightarrow ApplicationSemantics
\end{gathered}}
\]
它不是说所有真实网络都只有这七步，而是提示学习时分别去检查：主机网络配置是否已有；域名怎样解析；目的 IP 是同网段还是经网关；下一跳 MAC 如何得到；帧和分组怎样逐跳前进；目的主机怎样交给正确端口/连接；最后 bytes 怎样被 HTTP/DNS 等应用协议解释。

具体第一次访问可近似追：
\[
\begin{gathered}
\text{DHCP（若尚无配置）}\rightarrow DNS\rightarrow SameSubnet/DefaultGateway\rightarrow ARP\\
\downarrow\\
Frame/Switch\rightarrow Router/LPM\rightarrow TCP\rightarrow HTTP
\end{gathered}
\]
每一段只负责调用相应专题，绝不在总图里展开协议细节。

\section{网络最重要的概念边界：区别在哪里，题目怎么用}
\begin{longtable}{L{3.6cm}L{4.3cm}L{4.0cm}L{\dimexpr\linewidth-13.17cm\relax}}
\toprule
\textbf{概念边界} & \textbf{真正区别} & \textbf{题目信号} & \textbf{混淆后的错误}\\
\midrule
\key{Name} $\neq$ \key{IP} $\neq$ \key{MAC} $\neq$ \key{Port} & 名字指应用对象；IP 定位网络层目的；MAC 服务当前链路交付；Port 区分主机内端点 & DNS、子网/路由、ARP/Ethernet、socket/port 分别对应不同作用域 & 把“下一跳 MAC”当成最终目标主机 MAC，或认为 Port 用来路由 \\
\midrule
\key{Routing} vs. \key{Forwarding} & Routing 生成/更新规则；Forwarding 使用现有规则处理当前 packet & DV / LS / RIP / OSPF / BGP 是 routing；LPM 查表是 forwarding & 用 Dijkstra 解释每个 packet 的实时转发过程 \\
\midrule
\key{Reliability} $\neq$ \key{Flow} $\neq$ \key{Congestion} & 分别保护数据语义、接收方、网络路径 & seq/ACK；rwnd；cwnd/slow start 是三组信号 & 看到“窗口变小”就无法判断为什么变 \\
\midrule
\key{Transmission Delay} $\neq$ \key{Propagation Delay} & 一个取决于 packet 长度和发送速率 $L/R$；一个取决于距离和传播速度 $D/v$ & 改 packet 大小只影响 transmission；改距离只影响 propagation & 时延题把带宽和距离混进同一个公式 \\
\midrule
\key{Connection} $\neq$ \key{Physical Path} & TCP connection 是端点维护的逻辑状态；中间真实路径可由多条链路和路由器组成并可能变化 & 题目问 socket / 4-tuple 是 connection；问 router / link 是 path & 认为 TCP 建连时沿途路由器都保存连接状态 \\
\midrule
\key{Link Reliability} $\neq$ \key{End-to-End Reliability} & 都可用 ACK/重传，但一个修复一跳，一个保证两个端点之间的服务语义 & Wi-Fi / 链路重传与 TCP 重传 Scope 不同 & 因链路可靠就认为端到端绝不会丢 \\
\midrule
\key{TCP State} $\neq$ \key{Router State} & seq / window 属端点；prefix / next-hop 属路由器 & 看状态字段“谁需要读它做决定” & 把 TCP ACK 变化用于更新普通路由器转发表 \\
\midrule
\key{Layering} $\neq$ \key{Physical Law} & 分层是隔离变化的架构/分析方法，不是自然界强制边界 & NAT / VPN / Proxy 等现实实现会跨传统层次，但分析仍可分 Scope & 认为任何功能只能严格存在于某一层，无法理解跨层接口 \\
\bottomrule
\end{longtable}

\section{网络做题入口：八个词怎样变成动作}
\begin{examBox}{网络八步检查链}
\[
\boxed{\begin{gathered}
Scope\rightarrow Object\rightarrow StateOwner\rightarrow Event\\
\downarrow\\
Transition\rightarrow Feedback\rightarrow Bottleneck/Queue\rightarrow Target/Cost
\end{gathered}}
\]
\begin{enumerate}
  \item \key{Scope：}先说这是应用、端到端、跨网路径、一跳还是物理信道的问题。
  \item \key{Object：}当前处理的是 message、segment、datagram、frame 还是 bit/symbol。
  \item \key{StateOwner：}做决定所需的表/窗口/序号到底存在哪个端点或设备。
  \item \key{Event：}收到 packet/ACK/routing update，还是 timer/collision 等事件？
  \item \key{Transition：}哪张表、哪个窗口、哪个序号或队列发生改变？
  \item \key{Feedback：}这个动作会向谁发什么信息，怎样改变远端下一步判断？
  \item \key{Bottleneck/Queue：}如果是性能问题，哪个有限资源真正成为瓶颈，是否排队？
  \item \key{Target/Cost：}最后才计算时延、吞吐、利用率、窗口、地址范围等目标量。
\end{enumerate}
\end{examBox}

\begin{methodbox}{小例子：路由器收到一个 IP 分组，问输出接口}
Scope=Network/当前路由器的数据平面；Object=IP datagram；StateOwner=该路由器 forwarding table；Event=packet arrival；Transition/Action=对 DestinationIP 做 Longest Prefix Match，得到 NextHop/Interface；这里通常不需要分析 TCP seq，也不需要运行 OSPF。若题目进一步问“这张表怎样学出来”，才切换到 Routing 专题。
\end{methodbox}

\section{一页压缩}
\begin{corebox}{闭眼后应该剩下什么}
\[
\boxed{\text{Scope}\rightarrow\text{Local State}\xrightarrow{\text{Message/Timer}}\text{Transition}\rightarrow\text{Feedback}\rightarrow\text{Cost}}
\]
把它翻译成中文就是：\key{这个决定负责多远？谁手里有什么局部信息？收到什么或等了多久？它因此改了什么并做了什么？怎样通过消息影响别的节点？最后付出了什么代价？} 这才是网络母模型，而不是一串英文标签。
\end{corebox}

% ============================================================
\chapter{四科接口：真正的系统不是四块并排知识}
% ============================================================

\section{数据结构 × 计组：渐近复杂度为什么不是实际机器时间}
数组与链表可能都有线性遍历复杂度，但连续内存更容易利用 cache line 和预取；B/B+ 树的存在则直接说明：当随机 I/O 成为主成本时，结构目标从减少比较变成减少层数和 I/O。\key{算法成本必须最终落到底层机器成本模型。}

\section{计组 × OS：机制与策略的软硬件边界}
硬件提供特权级、MMU/TLB、trap/interrupt、timer、atomic instruction、DMA 等受控机制；OS 建立页表、调度策略、等待队列、文件和设备抽象，并决定何时使用这些机制。

\[
\boxed{\text{Hardware provides enforceable mechanisms; OS maintains abstractions, state and policy.}}
\]

\section{OS × 网络：packet 最终必须变成进程可见的数据}
网络协议不是漂浮在应用之下的抽象云。接收路径最终会经历 NIC 收帧、DMA/内存、内核协议栈、socket buffer、等待进程唤醒与调度，再回到用户态。因此 socket 是网络语义与 OS 资源/执行模型的关键接口。

\section{一条“应用读取网络数据”的跨科路径}
\arrowchain{远端 Application → Transport/IP/Link → NIC → DMA/Main Memory → Interrupt/Kernel → Network Stack → Socket Buffer → Blocked Process Wakeup → Scheduler → User read/recv → Application}

在这条路径里：
\begin{itemize}
  \item 网络回答“远端数据怎样抵达本机”；
  \item 计组回答“CPU、Cache、内存、DMA、中断怎样真实执行”；
  \item OS 回答“内核怎样组织 buffer、进程等待和控制权”；
  \item 数据结构回答“队列、散列、树、ring buffer 等状态怎样被高效组织”。
\end{itemize}

\section{跨科综合题的统一顺序}
\begin{examBox}{先分层，再追踪：箭头表示思考顺序}
\[
\boxed{\begin{gathered}\text{Object}\rightarrow\text{Layer/Scope}\rightarrow\text{State/Representation}\rightarrow\text{Event/Operation}\\\downarrow\\\text{Mapping/DataPath}\rightarrow\text{Queue/Competition}\rightarrow\text{Invariant}\rightarrow\text{Cost}\end{gathered}}
\]
\begin{enumerate}
  \item \key{Object：}当前到底在追 packet、instruction、process、page 还是数据结构节点？
  \item \key{Layer/Scope：}它现在处在网络哪一作用域、OS 哪个抽象、计组哪一级状态？先锁定世界。
  \item \key{State/Representation：}当前能直接查到哪些表、寄存器、队列、指针或缓存状态？
  \item \key{Event/Operation：}谁做了什么，导致状态有资格变化？
  \item \key{Mapping/DataPath：}需要经过哪条映射链或数据通路才能到下一层？
  \item \key{Queue/Competition：}有没有等待、冲突、带宽/端口/CPU 竞争？
  \item \key{Invariant：}每次状态变化后什么必须仍然正确？
  \item \key{Cost：}最终把完整路径换算成题目要求的时间、空间、I/O、吞吐或比较次数。
\end{enumerate}
\end{examBox}

以“应用阻塞读取网络数据”为例：Object 先从 remote packet 转为本机 socket/process；网络 Scope 负责分组抵达，本机 NIC/内存进入计组/OS 接口；State 包括 NIC descriptor、socket buffer 和 blocked process；packet arrival/DMA completion/interrupt 是事件；协议栈和 socket 映射是处理路径；进程等待队列体现 Queue；最终必须保证数据属于正确 socket 且用户缓冲区访问合法；最后才谈 DMA、唤醒、调度和 copy 的性能代价。这样一条综合链真正告诉你“何时切科”，而不是把四科名词堆在一起。

% ============================================================
\chapter{专题所有权与学习路线：总图怎样和后续手册配合}
% ============================================================

\section{Canonical Owner 原则}
同一知识点只能有一个专题负责完整解释。其他手册只能：
\begin{itemize}
  \item \key{Use：}调用结论；
  \item \key{Bridge：}解释接口；
  \item \key{Integration：}在真实生命周期中串联；
\end{itemize}
不能重新复制完整机制。

\begin{longtable}{L{3.7cm}L{4.3cm}L{\dimexpr\linewidth-9.27cm\relax}}
\toprule
\textbf{知识} & \textbf{Canonical Owner} & \textbf{典型调用者}\\
\midrule
Sliding Window / GBN / SR & NW-03 可靠传输 & TCP、网络性能综合\\
Flow Control & NW-06 Transport/TCP & TCP 综合\\
Congestion Control & NW-07 拥塞 & TCP/网络请求综合\\
ARP & NW-04 IP/Forwarding & NW-02 单跳、网络综合\\
Page Fault / COW & OS-04 VM & fork/mmap/OS 综合\\
DMA & OS-05 I/O + CO-07 硬件视角 & OS 综合、OS×Net\\
Open File State & OS-06 File System & fork/read/unlink 综合\\
Cache Mapping & CO-06 Cache & LOAD、性能综合\\
Graph Shortest Path & DS-07 Graph Algorithms & NW-05 Routing 的算法基础\\
\bottomrule
\end{longtable}

\section{推荐学习方式}
对任何专题使用同一节奏：
\[
\boxed{\text{总图定位}\rightarrow\text{专题打穿}\rightarrow\text{题目验证}\rightarrow\text{桥梁/综合重连}\rightarrow\text{回到一页压缩}}
\]
这条链不是出版流程，而是实际学习流程：
\begin{enumerate}
  \item \key{总图定位：}上课/看书前先知道这一章在母模型哪一格，例如“路由”主要深入 Local State + Message + Convergence；
  \item \key{专题打穿：}进入专题手册，把具体对象、算法、状态机、公式和题型学清；
  \item \key{题目验证：}做题时检查自己究竟卡在模型、识别、路径、执行还是检查；
  \item \key{桥梁/综合重连：}学完后把它接回邻近专题，例如 TCP 的可靠性调用 NW-03，而拥塞状态归 NW-07；
  \item \key{一页压缩：}最后重新用自己的话解释母模型，确认具体知识能重新挂回地图。
\end{enumerate}

\section{什么时候应该修改总手册}
只有出现以下变化才改总图：
\begin{enumerate}
  \item 某门课的母问题/母模型被证明不够准确；
  \item 专题边界或 canonical owner 发生变化；
  \item 出现能够解释多个专题的新上位关系；
  \item 长期做题证据证明某个关键概念边界必须进入学科总图。
\end{enumerate}
单道题的技巧、某个公式细节、某个协议字段、某种算法实现差异不直接进入总手册。

% ============================================================
\chapter{覆盖导航：总图只指出入口，不复制清单}
% ============================================================

详细考试覆盖清单统一由《408 四科统一方法论手册》维护；总图只保留四科覆盖域，防止同一清单在两本总册中分叉。

\begin{longtable}{L{2.2cm}L{5.2cm}L{\dimexpr\linewidth-8.67cm\relax}}
\toprule
\textbf{学科} & \textbf{覆盖域} & \textbf{检查入口}\\
\midrule
数据结构 & 复杂度、线性结构、树、图、查找与索引、串、内部与外部排序 & 先按核心模型定位 Relation、Operations、Representation、Invariant 与 Cost，再到四科方法论核对详细条目\\
计算机组成原理 & 数据表示、ISA、存储层次、数据通路与控制、流水线、总线与 I/O、性能 & 先定位软件可见语义、数据位置、时序与提交，再核对机制覆盖\\
操作系统 & 控制权、进程与调度、并发、虚拟内存、I/O、文件系统与持久化 & 先定位对象、关系、队列与事件，再进入对应 OS Topic\\
计算机网络 & 分层与性能、物理/链路、可靠传输、IP 与路由、TCP/UDP、应用层 & 先定位作用域和状态所有者，再进入对应网络专题\\
\bottomrule
\end{longtable}

\begin{boundarybox}{覆盖层的停止条件}
总图确认“专题入口存在”后即停止；具体公式、字段、算法步骤和题型清单只在四科方法论或专题手册维护。这样总图负责导航，深读手册负责解释，覆盖清单负责防漏。
\end{boundarybox}

% ============================================================
\chapter{最后压缩：考前真正应该留下什么}
% ============================================================

\begin{corebox}{四科一句话}
\begin{itemize}
  \item \key{数据结构：}为了操作需求，选择并维护合适的信息表示；
  \item \key{计组：}把 ISA 的软件语义实现成真实硬件的数据流与时间过程；
  \item \key{OS：}在硬件机制上维护受保护的对象、关系、队列和资源状态；
  \item \key{网络：}自治节点只凭局部状态和消息反馈共同建立端到端通信。
\end{itemize}
\end{corebox}

\begin{corebox}{四条母链}
\[
\text{DS: }Problem\rightarrow Relation\rightarrow Operations\rightarrow Representation\rightarrow Invariant\rightarrow Cost
\]
\[
\text{CO: }ISAState\rightarrow Data\rightarrow Path/Control\rightarrow Timing\rightarrow Commit\rightarrow Cost
\]
\[
\text{OS: }Objects/Relations/Queues\xrightarrow{Event+Mechanism+Policy}NewState\;|\;Invariant,Cost
\]
\[
\text{NET: }Scope\rightarrow LocalState\xrightarrow{Message/Timer}Transition\rightarrow Feedback\rightarrow Cost
\]
\end{corebox}

\begin{methodbox}{最终学习原则}
不要把“知道名词”误认为“拥有模型”，也不要把“会背母链”误认为“会用母模型”。真正掌握一个知识点，应当能够：
\begin{enumerate}
  \item 用自己的中文解释母模型每一格是什么意思；
  \item 指出它在教材/专题中的具体章节实例；
  \item 给一个小题，按箭头顺序真正走一遍；
  \item 遇到 $A\neq B$ 时说出判据、题目信号和混淆后果；
  \item 最终回答它解决什么问题、依赖什么状态/表示、由什么操作或事件触发、维护什么不变量、代价从哪里来、应该挂到哪个专题。
\end{enumerate}
\end{methodbox}

\begin{center}
\Large\bfseries\color{deep}
总图负责让你永远知道自己站在哪里；专题负责让你知道这一块为什么这样运行；题目负责证明这个模型是否真的属于你。
\end{center}

\end{document}
